\chapter{A Quick Primer On Compiler Verification}
\label{ch:background}
\epigraph{A description of Zaira as it is today should contain all Zaira’s
past. The city, however, does not tell its past, but contains it like the lines
of a hand, written in the corners of the streets, the gratings of the windows,
the banisters of the steps, the antennae of the lightning rods, the poles of
the flags, every segment marked in turn with scratches, indentations,
scrolls}{Italo Calvino, Invisible Cities}


% Two further mechanisms matter for this thesis. First, \emph{proof by
% reflection}~\citep{gonthier2010introduction}: the typeclass
% \texttt{Decidable} associates to a proposition a program that decides it, and
% the \texttt{decide} tactic proves a proposition by running that program and
% appealing to its soundness proof. Reflection replaces a potentially enormous
% proof term by a computation. When kernel reduction is too slow for that
% computation, Lean offers the \texttt{ofReduceBool} axiom: the decision
% procedure is compiled to native code and its Boolean result is trusted. This
% extends the trusted code base from the kernel to the Lean compiler---a
% trade-off we revisit when discussing Lean's verified bitblaster \bvdecide{}
% (\cref{ch:background-bv}).
% Second, Lean's runtime uses
% reference counting and implements the \emph{functional but in-place} (FBIP)
% optimization~\citep{Ullrich2019}: a purely functional update to an object
% whose reference count is $1$ is performed by in-place mutation. Arrays and
% hash maps (which are built on arrays) therefore retain their usual imperative
% performance characteristics inside pure code, provided only a single reference
% to them is kept live.

\section{The SSA Compiler Intermediate Representation}
\label{sec:bg:ssa}

The bitvector problems this thesis attacks arise from compiler intermediate
representations, so we recall the relevant vocabulary. Static single
assignment (SSA) form~\citep{rastello2022ssa} is the dominant style of
compiler IR: every variable is assigned exactly once, which makes def-use
relationships immediate and thereby simplifies analyses and transformations.
\emph{Peephole optimizations}~\citep{mckeeman1965peephole}, which replace
short instruction sequences by semantically equivalent cheaper ones, are the
workhorse transformation on SSA IRs; roughly 10\% of all IR-transforming code
in LLVM~\citep{lattner2004llvm} belongs to its InstCombine peephole
optimizer. Crucially, peephole rewriting on SSA operates along the
\emph{def-use chain} rather than on syntactically contiguous instructions: the
rewrite
$(y = x + 1;\ z = y - 1) \mapsto (z = x)$ applies to any program containing
$(y = x + 1;\ \ldots;\ z = y - 1)$, regardless of the interleaved
instructions---unlike the classical straight-line peephole rewriting over
assembly mechanized by Peek~\citep{mullen2016verified}.

The MLIR compiler framework~\citep{lattner2021mlir} lowers the cost of
instantiating \emph{domain-specific} SSA IRs. An MLIR \emph{dialect} packages
a set of operations and types at a chosen abstraction level, and dialects for
machine learning, quantum computing, and fully homomorphic encryption
(FHE)~\citep{gentry2009fully,viand2022heco} coexist in one infrastructure.
High-level dialects often have \emph{value semantics} (referential
transparency), enabling local, side-effect-free algebraic reasoning that is
simply unavailable at the level of LLVM-style load/store/branch code: an FHE
rewrite such as $a + X^{2^n} + 1 \mapsto a$ is a one-line quotient-ring
identity in a polynomial dialect, but a heroic exercise in reasoning about
loops and memory once lowered to LLVM. MLIR additionally introduces
\emph{regions}, which let operations contain nested code and so express
structured control flow (loops, conditionals) without unstructured branches,
simplifying termination and semantics arguments.

Building on this vocabulary, we develop Lean-MLIR~\citep{leanmlir}, a
framework for intrinsically well-typed SSA with regions in Lean, generic over
a user-defined dialect, with a verified peephole rewriter whose local rewrite
proofs lift to whole-program semantic preservation, and proof automation that
reduces each rewrite obligation to a clean mathematical goal
(\cref{ch:lean-mlir}).

\section{The MLIR Compiler Framework}

\section{Approaches to Verifying MLIR-Style Compilers}
\label{sec:bg:compiler-verification}

How does one come to trust a compiler? The most ambitious answer is to verify
it once and for all: write the compiler inside a proof assistant and prove,
pass by pass, that every translation preserves the semantics of its input.
CompCert~\citep{leroy2009formally,leroy2016compcert} is the landmark result of
this school---a moderately optimizing C compiler, programmed and proved in
Coq, whose correctness theorem stretches from the parsed C program down to
assembly---and CakeML~\citep{kumar2014cakeml} plays the same role for an ML
dialect in HOL4. The guarantee is as strong as guarantees get; the price is
proof engineering that must move in lockstep with the compiler. Even pinning
down the semantics of an industrial IR is a research programme of its own:
Vellvm~\citep{zhao2012formalizing,zhao2013formal} and K-LLVM~\citep{li2020k}
each mechanize substantial fragments of LLVM, and neither covers the compiler
that consumes it. Production compilers, meanwhile, grow by thousands of
commits a month; verified reimplementations cannot chase them.

\emph{Translation validation}~\citep{pnueli1998translation} lowers the
ambition and, in exchange, keeps pace: instead of verifying the compiler,
verify each of its runs, by checking after the fact that the emitted program
refines the source program. The compiler remains an untrusted black box; only
the checker need be sound, and the checker only ever looks at pairs of
concrete programs.

The Alive line of tools brought this style, powered by SMT solvers, into the
day-to-day of the LLVM community. Alive~\citep{lopes2015provably} offers a
DSL in which an InstCombine rewrite is stated once and discharged
automatically---or refuted with a counterexample---and
Alive2~\citep{lopes2021alive2} validates runs of the existing LLVM
optimizer wholesale; it is the most referenced tool in LLVM's commit log.
Alive-tv~\citep{DBLP:conf/cav/BangNCJL22} extends the approach to MLIR's
tensor dialects. Push-button convenience, however, has a cost in trust: the
SMT solver, and the encoding of IR semantics into SMT, both join the trusted
base. AliveInLean~\citep{lee2019aliveinlean} verifies the encoding, but the
solver remains an oracle. Moreover, the theory being decided is that of
\emph{fixed-width} bitvectors, while an InstCombine rewrite must be correct at
every register width; Alive checks a handful of widths and hopes the rest
behave alike.

Between these poles sits \emph{mechanized rewriting}: prove a rewriting
\emph{engine} correct once, inside a proof assistant, and then prove each
rewrite against the real semantics of the IR, with no solver and no width
bound in the trusted base. Peek~\citep{mullen2016verified} did this for
assembly peepholes inside CompCert. The bet this thesis makes is that the two
ends of the spectrum can be joined: an interactive theorem prover supplies the
trusted foundation and the rewriting framework (\cref{ch:lean-mlir}), and
verified decision procedures supply the push-button automation that makes the
framework usable at scale
(\cref{ch:mono-width,ch:multi-width,ch:floating-point}). To state rewrites in
the first place, we need the vocabulary of the IRs they live in, which we now
recall.


\section{The Lean Proof Assistant}

In this thesis, we mechanize our results in the Lean proof assistant.
Lean~\citep{demoura2021lean} is a dependently typed interactive theorem
prover (ITP) that doubles as a general-purpose, purely functional programming
language. We mechanize our compiler verification framework,
as well as the decision procedures that power it, in Lean.
